El aseguramiento de la seguridad operacional constituye el mecanismo mediante el cual SKYFILMS verifica de manera sistemática la eficacia de su Sistema de Gestión de la Seguridad Operacional (SMS). Esta sección establece los procesos que permiten evaluar el desempeño de la organización frente a sus objetivos de seguridad, garantizar la gestión controlada de los cambios y promover la mejora continua.\\

El proceso general de gestión del rendimiento en seguridad operacional y su relación con los sistemas de recopilación y procesamiento de datos sobre seguridad operacional (SDCPS) se resume en la siguiente tabla:

\begin{table}[H]
\centering
\renewcommand{\arraystretch}{1.35}
\begin{tabularx}{\textwidth}{|>{\RaggedRight\arraybackslash}p{4.6cm}|>{\RaggedRight\arraybackslash}X|}
\hline
\rowcolor[gray]{0.9}\textbf{Elemento del SDCPS} & \textbf{Qué hace} \\ \hline
\rowcolor[HTML]{FFE04D}\textbf{Rendimiento del SMS} & Es el centro del sistema: todos los demás elementos lo alimentan y se retroalimentan con él. \\ \hline
Definición y actualización de SPI & Define los indicadores de seguridad operacional (Anexo 9, tabla A) y los actualiza. \\ \hline
Definición y actualización de KPI & Define los indicadores internos de gestión (Anexo 9, tabla B). \\ \hline
Medición del rendimiento & Mide cada mes los ciclos de vuelo y los eventos y calcula las tasas. \\ \hline
Análisis del SMS & El SAG analiza las tendencias, causas y alertas. \\ \hline
Promoción del SMS & Comunica los resultados y las lecciones aprendidas al personal. \\ \hline
Identificación de medidas necesarias & Define las acciones y planes de mejora que se derivan de los resultados. \\ \hline
\end{tabularx}
\caption{Sistemas de recopilación y procesamiento de datos sobre seguridad operacional (SDCPS).}
\end{table}


\section{Observación y Medición del Rendimiento}
SKYFILMS implementará un sistema estructurado de observación y medición del rendimiento en seguridad operacional, en cumplimiento del RAC 219, la Circular MAUT-1.0-22-005 de Aerocivil y las directrices del Doc 9859 de la OACI.

\begin{itemize}
    \item La organización establecerá y mantendrá Indicadores de Desempeño en Seguridad Operacional (SPI), los cuales estarán directamente vinculados con los objetivos de seguridad de SKYFILMS y con los niveles aceptables de seguridad operacional aprobados por la Autoridad Aeronáutica.
    \item La recopilación, análisis y reporte de los SPI se realizará de forma periódica, garantizando la trazabilidad de la información y permitiendo identificar tendencias, desviaciones y oportunidades de mejora.
    \item Los resultados serán evaluados anualmente, y servirán de base para la toma de decisiones en materia de mitigación de riesgos, asignación de recursos y definición de acciones correctivas y preventivas.
    \item Toda la información asociada al monitoreo de SPI se documentará en formatos oficiales de la organización y será consolidada en un anexo específico al presente manual \textbf{(Anexo 9 - Indicadores de Rendimiento de Seguridad Operacional SPI's)}, que incluirá la descripción detallada de cada indicador, su fórmula, metas, plan de acción y mecanismos de evaluación.
\end{itemize}



Los SPI de Skyfilms son los once indicadores que la Aerocivil concertó para explotadores UAS (circular MAUT-1.0-22-005). Sus fórmulas (eventos por cada 1\,000 ciclos de vuelo), metas, niveles de alerta y forma de entrega están en el Anexo 9, y su relación con los objetivos de seguridad operacional en el Anexo 1. El Gerente SMS anota los datos cada mes, los presenta al SAG cada trimestre y el SRB los revisa al menos una vez al año. Cada año, antes del 30 de marzo, se envía el archivo SPI a la Aerocivil.

\section{Gestión del Cambio}
SKYFILMS reconoce que la introducción de cambios en el funcionamiento de la organización puede generar nuevos riesgos o modificar los existentes, debido a esto a continuación se listan los principales desencadenantes que activan el proceso de gestión del cambio:

\begin{enumerate}
    \item Cambios las operaciones.
    \item Cambios o inclusión de equipos.
    \item Cambios o inclusión de personal.
    \item Cambios o inclusión de procedimientos.
\end{enumerate}

 Por lo tanto, se aplicará un proceso formal de gestión del cambio \textbf{(Anexo 10 - Procedimiento de Gestión del Cambio)} para garantizar que toda modificación sea evaluada y aprobada antes de su implementación.


\textbf{Desencadenantes, responsables, plazos y seguimiento:} están en el Anexo 10. Incluyen la participación de las personas afectadas, la aprobación antes de aplicar el cambio, la revisión de las evaluaciones de riesgo ya hechas, el análisis de riesgos transitorios y acumulados, la comunicación y el seguimiento a los 30 y 90 días.

\section{Programa de Auditoria Interna}

\subsection{Propósito}

SKYFILMS establecerá y mantendrá un Programa de Auditoría Interna como herramienta sistemática, independiente y documentada para evaluar:

\begin{itemize}
    \item La conformidad del Sistema de Gestión de la Seguridad Operacional (SMS) con la normativa aplicable.

    \item El cumplimiento de las políticas, procedimientos y controles establecidos en el presente Manual.

    \item La eficacia de los controles de mitigación implementados.

    \item El desempeño de los Indicadores de Seguridad Operacional (SPI).

    \item La madurez y efectividad global del SMS.
\end{itemize}

El programa constituye un elemento fundamental del proceso de aseguramiento de la seguridad operacional y soporte a la mejora continua.

\subsection{Alcance}
El Programa de Auditoría Interna aplicará a:

\begin{itemize}
    \item Operaciones UAS.

    \item Gestión técnica y mantenimiento.

    \item Gestión documental del SMS.

    \item Gestión de riesgos.

    \item Gestión del cambio.

    \item Gestión de reportes voluntarios.

    \item Cumplimiento normativo (RAC 100, RAC 219 y demás regulación aplicable).

    \item Gestión de instrucción y competencia del personal.

    \item Cumplimiento de procedimientos operacionales.
\end{itemize}

\subsection{Periodicidad}
El programa anual de auditorías será aprobado por el Gerente SMS y el Gerente General, de la cual realizará como mínimo una auditoría interna \textbf{anual} completa del Sistema de Gestión de Seguridad Operacional SMS, y ademas podrán programarse auditorías parciales o específicas cuando:

\begin{itemize}
    \item Se identifiquen desviaciones significativas en los SPI.

    \item Existan hallazgos recurrentes.

    \item Se implementen cambios operacionales relevantes.

    \item La Dirección lo considere necesario.
\end{itemize}

\subsection{Metodologia}
Las auditorías internas se desarrollarán bajo las siguientes etapas:

\begin{enumerate}
    \item Planificación de la auditoría.

    \item Revisión documental.

    \item Entrevistas y verificación en campo.

    \item Evaluación de cumplimiento.

    \item Identificación de hallazgos.

    \item Emisión de informe formal.

    \item Seguimiento a acciones correctivas.

    \item Verificación de eficacia.
\end{enumerate}

\subsection{Calendario, Independencia y Competencia del Auditor}
\begin{itemize}
    \item \textbf{Calendario:} el programa anual se fija en enero. Se audita una vez al año cada área (operaciones, mantenimiento, documentos, riesgos, capacitación) y el SMS completo. Las fechas se registran en el programa y se avisan con 15 días de anticipación.
    \item \textbf{Independencia:} nadie audita su propio trabajo. El Gerente SMS audita operaciones y mantenimiento; al SMS lo audita el Jefe de Pilotos con apoyo de un auditor externo o de una persona sin relación con el área auditada.
    \item \textbf{Competencia del auditor:} curso de SMS y conocimiento del RAC 100 y el RAC 219; el curso de personal clave del Anexo 12 cumple este requisito.
    \item \textbf{Contratistas y proveedores:} se audita al proveedor de mantenimiento (cuando exista) y la ejecución de sus controles, al menos una vez al año.
    \item \textbf{Acceso al Ejecutivo Responsable:} el responsable de la auditoría entrega el informe directamente al Gerente General y puede pedirle una reunión en cualquier momento.
    \item \textbf{Informes y cierre:} el informe sale dentro de 15 días hábiles después de la auditoría; los hallazgos se cierran según el Plan de Mejora Continua (Anexo 11) y los registros se guardan cinco años.
\end{itemize}

\subsection{Clasificación de Hallazgos}

Los hallazgos derivados de auditoría se clasificarán como:\\

\textbf{Hallazgo Mayor:}

\begin{itemize}
    \item Incumplimiento que afecta directamente la seguridad operacional.

    \item Ausencia o falla crítica de un control.

    \item Incumplimiento normativo significativo.
\end{itemize}

\textbf{Hallazgo Menor:}

\begin{itemize}
    \item Desviación puntual sin impacto inmediato en la seguridad.

    \item Oportunidad de mejora.
\end{itemize}

\textbf{Observación:}

\begin{itemize}
    \item Recomendación preventiva sin incumplimiento directo.
\end{itemize}

\subsection{Gestión de Hallazgos}
Todos los hallazgos deberán:

\begin{itemize}
    \item Registrarse en el Formato de Plan de Mejora Continua.

    \item Definir acción correctiva o preventiva.

    \item Asignar responsable.

    \item Establecer plazo de cierre.

    \item Verificar eficacia antes de cierre formal.
\end{itemize}

El Gerente SMS realizará seguimiento periódico hasta su cierre.

\subsection{Informe de Auditoria}
Cada auditoría generará un informe formal que contendrá:

\begin{itemize}
    \item Alcance

    \item Objetivo

    \item Criterios utilizados

    \item Metodología aplicada

    \item Hallazgos identificados

    \item Nivel de cumplimiento

    \item Conclusión general sobre el estado del SMS
\end{itemize}

El informe será presentado a la Gerencia para su conocimiento y aprobación de acciones.

\subsection{Conservación de Registros}
Todos los registros asociados al Programa de Auditoría Interna serán conservados conforme al sistema de gestión documental del SMS por un \textbf{período mínimo de cinco (5) años}.

\subsection{Integración con el Sistema SMS}

El Programa de Auditoría Interna alimenta directamente:

\begin{enumerate}
    \item La medición del desempeño del SMS.

    \item La revisión de los SPI.

    \item La gestión del cambio.

    \item El Plan de Mejora Continua.

\end{enumerate}

Lo que constituye un mecanismo estructural para asegurar la efectividad y el seguimiento de los datos utilizados para la toma de decisiones de gestión y asegurar la mejora continua sostenida del Sistema de Gestión de Seguridad Operacional de SKYFILMS.

\section{Mejora Continua}
La seguridad operacional es un proceso dinámico que requiere de revisión y optimización permanentes. SKYFILMS se compromete a mantener un ciclo de mejora continua del SMS, garantizando que los resultados obtenidos conduzcan a un fortalecimiento progresivo del sistema.

\begin{itemize}
    \item Se realizarán auditorías internas de seguridad operacional con una periodicidad mínima \textbf{anual}, así como revisiones por la dirección al menos una vez por año.
    \item Los hallazgos derivados de auditorías, inspecciones, encuestas de seguridad, reportes voluntarios y el análisis de los SPI serán integrados en un Plan de Mejora Continua \textbf{(Anexo 11 - Plan de Mejora Continua)}, que definirá acciones correctivas, responsables, plazos y mecanismos de verificación.
    \item Toda acción correctiva o preventiva deberá estar sustentada en datos objetivos y será verificada para confirmar su eficacia antes de ser cerrada.
    \item La evolución del desempeño de seguridad operacional, las lecciones aprendidas y las mejoras implementadas serán comunicadas de manera transparente a todo el personal, fortaleciendo la cultura de seguridad en la organización.
    \item Como responsable principal esta el Gerente SMS quien sera el encargado de liderar las auditorias y revisiones.
\end{itemize}

